\documentclass{article}
\usepackage{graphicx}
\usepackage{amsmath}
\usepackage{caption}
\usepackage{subcaption}
\usepackage{geometry}
\geometry{margin=1in}

\title{Assignment 3 - CIFAR-10 Classification with Convolutional Neural Networks}
\author{Martín Bravo, DD2424 Deep Learning, KTH Royal Institute of Technology}
\date{\today}

\begin{document}

\maketitle


\section{Introduction}
This assignment explores the implementation and optimization of a convolutional  neural network for image classification on the CIFAR-10 dataset. We investigate several key aspects of training deep learning models, including the architecture of the network, learning rate schedules, and regularization techniques. The goal is to achieve high classification accuracy while maintaining efficient training times.

In this assignment, we explored the implementation and optimization of a convolutional neural network for image classification on the CIFAR-10 dataset. We investigated various techniques, including varying the architecture, using cyclic learning rates, and applying label smoothing.

Through our experiments, we found that increasing the cyclic learning rate can improve training performance, but it does not always lead to better generalization. We also discovered that label smoothing is an effective regularization technique that helps to reduce overfitting and improve the model's generalization capabilities. Overall, this assignment provided valuable insights into the design and optimization of convolutional neural networks for image classification tasks.


\section{Dataset and Preprocessing}
The CIFAR-10 dataset consists of $32\times32$ color images belonging to 10 classes. We normalized the input data per feature using the training set mean and standard deviation. The dataset was split into training, validation, and test sets. Below are a few CIFAR-10 examples:

\begin{figure}[h!]
    \centering
    \includegraphics[width=0.6\textwidth]{imgs/assignment1_cifar_examples.png}
    \caption{Sample CIFAR-10 images from the dataset}
\end{figure}

\section{Background}

\subsection{Convolutional Neural Network Architecture}
The architecture of the network is a 3-layer network. The first layer is a convolutional layer with ReLU activation, followed by a fully connected layer with ReLU activation, and finally an output layer with softmax activation.
The input to the network is a flattened image vector $x$ of size $3072 \times 1$ (since each CIFAR-10 image has dimensions $32 \times 32 \times 3$). The parameters of the network include:

\begin{itemize}
    \item A convolutional layer with $n_f$ filters of size $f \times f \times 3$, stride $f$.
    \item A weight matrix $W_1$ of size $n_p n_f \times m$, where $n_p = (32/f) \times (32/f)$ is the number of patches after convolution, and $m$ is the number of hidden units in the fully connected layer.
    \item A bias vector $b_1$ of size $m \times 1$.
    \item A weight matrix $W_2$ of size $m \times 10$, where $10$ is the number of classes.
    \item A bias vector $b_2$ of size $10 \times 1$.
\end{itemize}

To compute the output of the network, we first compute the scores for each class as:
\[h = \text{Conv}(x, W_{conv})\]
where $\text{Conv}$ denotes the convolution operation with the filters $W_{conv}$. The output $h$ is then flattened into a vector of size $n_p n_f \times 1$ and passed through the fully connected layers:
\[s_1 = W_1 h + b_1, \quad z = \text{ReLU}(s_1), \quad s_2 = W_2 z + b_2, \quad p = \text{Softmax}(s_2)\]

The predicted class is then given by:
\[
k^* = \arg\max_{1 \leq k \leq 10} p_k
\]

This allows more complexity compared to the simple linear classifier, as the hidden layer allows for non-linear transformations of the input.

\subsection{Cyclic Learning Rate}

The cyclic learning rate is a technique that varies the learning rate during training. It starts at a low value, increases linearly to a high value, and then decreases linearly back to the low value. This is done in a cyclic manner, hence the name.

We use a linear increase and decrease of the learning rate, with a period of $T$ iterations, and a step size of $n_s$. The learning rate is given by:

\begin{equation}
    \eta(t) = \eta_{min} + \frac{t}{n_s} (\eta_{max} - \eta_{min})
\end{equation}

where $t$ is the current iteration, $T$ is the total number of iterations, $n_s$ is the step size, $\eta_{min}$ is the minimum learning rate, and $\eta_{max}$ is the maximum learning rate.


\subsection{Increasing Learning Rate} 

Another technique we explored is increasing the learning rate over time. This is done by starting with a step size of $n_s$ and increasing it by a factor of 2 every cycle. This allows the model to be trained for longer periods.

\subsection{Label Smoothing}

Label smoothing is a regularization technique that replaces the one-hot encoded labels with a smoothed version. Instead of using a hard label of 1 for the correct class and 0 for the incorrect classes, we use a soft label that assigns a small probability to the incorrect classes. This helps to prevent overfitting and improves generalization.

\subsection{Checking gradients computation}

In order to check the correctness of the gradients computation, we compared our implementation with the one provided by PyTorch. To measure how close the gradients are, we computed the relative error of the gradients:

\begin{equation}
    \epsilon = \frac{\| \nabla J - \nabla J_{\text{PyTorch}} \|}{\| \nabla J \| + \| \nabla J_{\text{PyTorch}} \|} \leq \delta
\end{equation}

We test that the distance between the gradients is less than a constant value $\delta$.

\section{Results}
We conducted three main experiments, first we varied the number of filters and stride in the convolutional layer, then we explored the effect of using an increasing cyclic learning rate, and finally we tested the impact of label smoothing on the model's performance.


\subsection{Experiment 1: Varying number of filters and stride}

The convolutional layer's parameters, specifically the number of filters ($n_f$) and the stride ($f$), play a crucial role in determining the model's capacity and performance. We experimented with different configurations to observe their effects on training and validation accuracy.

\begin{figure}[h!]
    \centering
    \begin{subfigure}{0.24\textwidth}
        \includegraphics[width=\textwidth]{imgs/accuracy_arch_arch1.pdf}
        \caption{$f$=2, $n_f$=3}
    \end{subfigure}
    \hfill
    \begin{subfigure}{0.24\textwidth}
        \includegraphics[width=\textwidth]{imgs/accuracy_arch_arch2.pdf}
        \caption{$f$=4, $n_f$=10}
    \end{subfigure}
    \hfill
    \begin{subfigure}{0.24\textwidth}
        \includegraphics[width=\textwidth]{imgs/accuracy_arch_arch3.pdf}
        \caption{$f$=8, $n_f$=40}
    \end{subfigure}
    \hfill
    \begin{subfigure}{0.24\textwidth}
        \includegraphics[width=\textwidth]{imgs/accuracy_arch_arch4.pdf}
        \caption{$f$=16, $n_f$=160}
    \end{subfigure}
    \vfill
    \begin{subfigure}{0.24\textwidth}
        \includegraphics[width=\textwidth]{imgs/loss_arch_arch1.pdf}
        \caption{$f$=2, $n_f$=3}
    \end{subfigure}
    \hfill
    \begin{subfigure}{0.24\textwidth}
        \includegraphics[width=\textwidth]{imgs/loss_arch_arch2.pdf}
        \caption{$f$=4, $n_f$=10}
    \end{subfigure}
    \hfill
    \begin{subfigure}{0.24\textwidth}
        \includegraphics[width=\textwidth]{imgs/loss_arch_arch3.pdf}
        \caption{$f$=8, $n_f$=40}
    \end{subfigure}
    \hfill
    \begin{subfigure}{0.24\textwidth}
        \includegraphics[width=\textwidth]{imgs/loss_arch_arch4.pdf}
        \caption{$f$=16, $n_f$=160}
    \end{subfigure}
    \caption{Training and validation accuracy and loss for different architectures}
    \label{fig:arch1_4}
\end{figure}

\begin{figure}[h!]
    \centering
    \includegraphics[width=0.3\textwidth]{imgs/compare_times.pdf}
    \caption{Comparison of training times for different architectures}
    \label{fig:compare_times}
\end{figure}

We can see that adding the convolutional layer, even with a small number of filters and a small stride, such as $f$=2, $n_f$=3, significantly improves the model's performance compared to a simple linear classifier. As we increase the number of filters and the stride, the model's capacity increases, leading to better training and validation accuracy. However, we also observe that larger architectures take way more time to train, as shown in Figure \ref{fig:compare_times}. Besides, very large architectures tend to overfit the training data, leading to a huge gap between training and validation accuracy as we can see in the case of $f$=16, $n_f$=160. We found that the best trade-off between performance and training time was achieved with $f$=4, $n_f$=10 and $f$=8, $n_f$=40, achieving validation accuracies of around 60\% in a reasonable training time and without overfitting too much. We will use these architectures for the next experiments.

\subsection{Experiment 2: Increasing Cyclic Learning Rate}

We implemented an increasing cyclic learning rate, where the step size $n_s$ is doubled every cycle. This allows the model to be trained for longer periods, potentially leading to better convergence and improved performance. We tested this technique on the two best architectures from the previous experiment: $f$=4, $n_f$=10 and $f$=8, $n_f$=40. Additionally, we tested the effect of increasing the number of filters to $n_f$=40 for the $f$=4 architecture.

\begin{figure}[h!]
    \centering
    \begin{subfigure}{0.32\textwidth}
        \includegraphics[width=\textwidth]{imgs/accuracy_arch_arch2_long.pdf}
        \caption{$f$=4, $n_f$=10, increased cyclic LR}
    \end{subfigure}
    \hfill
    \begin{subfigure}{0.32\textwidth}
        \includegraphics[width=\textwidth]{imgs/accuracy_arch_arch3_long.pdf}
        \caption{$f$=8, $n_f$=40, increased cyclic LR}
    \end{subfigure}
    \hfill
    \begin{subfigure}{0.32\textwidth}
        \includegraphics[width=\textwidth]{imgs/accuracy_arch_arch2_long_more_filters.pdf}
        \caption{$f$=4, $n_f$=40, increased cyclic LR}
    \end{subfigure}
    \vfill
    \begin{subfigure}{0.32\textwidth}
        \includegraphics[width=\textwidth]{imgs/loss_arch_arch2_long.pdf}
        \caption{$f$=4, $n_f$=10, increased cyclic LR}
    \end{subfigure}
    \hfill
    \begin{subfigure}{0.32\textwidth}
        \includegraphics[width=\textwidth]{imgs/loss_arch_arch3_long.pdf}
        \caption{$f$=8, $n_f$=40, increased cyclic LR}
    \end{subfigure}
    \hfill
    \begin{subfigure}{0.32\textwidth}
        \includegraphics[width=\textwidth]{imgs/loss_arch_arch2_long_more_filters.pdf}
        \caption{$f$=4, $n_f$=40, increased cyclic LR}
    \end{subfigure}
    \caption{Training and validation accuracy and loss for architectures with increased cyclic learning rate}
    \label{fig:arch2_3_long}
\end{figure}

We can see that using an increasing cyclic learning rate leads to improved training accuracy. Almost all architectures reach training accuracies above 80\%. However, the validation accuracy does not improve significantly, where most of the architectures reach around 60\%. This indicates that while the model is able to fit the training data better, it does not generalize much better to unseen data. 
The architecture with $f$=4 and $n_f$=10 reaches a validation accuracy of around 60\%, but having a good ratio of training to validation accuracy, suggesting that it generalizes better than the other architectures.
The architecture with $f$=8 and $n_f$=40 achieves the best validation accuracy of around 62\%, suggesting that increasing the number of filters can help improve performance, but at the cost of having overfitting more to the training data (training accuracy of almost 90\%).
Finally, the architecture with $f$=4 and $n_f$=40 achieves a similar problem, the validation accuracy increases but the training accuracy is much higher, indicating overfitting. We can conclude that while increasing the cyclic learning rate helps improve training performance, it does not necessarily lead to better generalization, and care must be taken to avoid overfitting. We will tackle this in the next experiment using label smoothing.

\subsection{Experiment 3: Label Smoothing}

We implemented label smoothing as a regularization technique to help improve generalization and reduce overfitting. We tested this technique on the architecture with $f$=4, $n_f$=40 and $n_h$=300, which showed signs of overfitting in the previous experiment. We also use cyclic learning rates with an extra cycle of training:

\begin{figure}[h!]
    \centering
    \begin{subfigure}{0.49\textwidth}
        \includegraphics[width=\textwidth]{imgs/accuracy_arch_arch5_baseline.pdf}
        \caption{$f$=4, $n_f$=40, $n_h$=300, baseline}
    \end{subfigure}
    \hfill
    \begin{subfigure}{0.49\textwidth}
        \includegraphics[width=\textwidth]{imgs/accuracy_arch_arch5_label_smoothing.pdf}
        \caption{$f$=4, $n_f$=40, $n_h$=300, label smoothing}
    \end{subfigure}
    \vfill
    \begin{subfigure}{0.49\textwidth}
        \includegraphics[width=\textwidth]{imgs/loss_arch_arch5_baseline.pdf}
        \caption{$f$=4, $n_f$=40, $n_h$=300, baseline}
    \end{subfigure}
    \hfill
    \begin{subfigure}{0.49\textwidth}
        \includegraphics[width=\textwidth]{imgs/loss_arch_arch5_label_smoothing.pdf}
        \caption{$f$=4, $n_f$=40, $n_h$=300, label smoothing}
    \end{subfigure}
    \caption{Training and validation accuracy for architecture 5 with and without label smoothing}
    \label{fig:arch5}
\end{figure}

We can see that using label smoothing helps to reduce the previous problem, as the gap between training and validation loss is smaller compared to the baseline. Analyzing the loss plots, label smoothing helps to stabilize the training process, leading to a smoother decrease of the gap between training and validation loss. Overall, label smoothing proves to be an effective regularization technique that enhances the model's analysis step and generalization capabilities.

\section{How to boost the performance even more}

To further boost the performance of the convolutional neural network on the CIFAR-10 dataset, several strategies can be employed, one of the main problems here was the overfitting of the model. Varying the lambda term would have helped to reduce overfitting even more. Additionally, we could experiment with more advanced architectures, such as deeper networks (adding more convolutional and fully connected layers), imitating architectures like ResNet.

\section{Conclusion}

In this assignment, we explored the implementation and optimization of a convolutional neural network for image classification on the CIFAR-10 dataset. We investigated various techniques, including varying the architecture, using cyclic learning rates, and applying label smoothing.

Through our experiments, we found that increasing the cyclic learning rate can improve training performance, but it does not always lead to better generalization. We also discovered that label smoothing is an effective regularization technique that helps to reduce overfitting and improve the model's generalization capabilities. Overall, this assignment provided valuable insights into the design and optimization of convolutional neural networks for image classification tasks.


\end{document}

